% ctx-texto.tex
%
% Copyright (C) 2022--2026 José A. Navarro Ramón <janr.devel@gmail.com>
% Licencia del código GPLv2
% Licencia Creative Commons Recognition Non-Commercial Share-alike.
% (CC-BY-NC-SA)


\startcomponent ctx-texto
  \product tutorialLuaMetaFun
  
  \startchapter[title=Texto]
    En este capítulo veremos cómo utilizar texto en figuras realizadas con
    \MetaFun/\MetaPost.

    La primera línea del código que sigue se analizó en el capítulo anterior.
    Analizaremos la segunda línea y su resultado en \MetaFun:
    \starttyping[option=METAFUN]
      \startMPcode
        fill fullcircle xyscaled (8cm,1cm) withcolor "darkred" ;
        draw textex{"\bf Texto A"} withcolor "white" ;
      \stopMPcode
    \stoptyping
  
    \startMPcode
      fill fullcircle xyscaled (8cm,1cm) withcolor "darkred" ;
      draw textext("\bf This is text A") withcolor "white" ;
    \stopMPcode

    En esta segunda línea utilizamos el comando \type{draw}, que dibujará el trazo que
    indiquemos a continuación. La que sigue es la función generadora \type{textext}.
  
    \startsection[title=Función generadora \color[red]{\type{textext}}]
    
      Esta es una macro específica de \ConTeXt / \MetaFun\ (no existe en \MetaPost\
      clásico estándar).
      Técnicamente, es una función generadora que toma una cadena de texto, la procesa
      a través de \ConTeXt\ y devuelve un objeto gráfico de tipo \type{picture},
      que es el sujeto o el objeto que vas a pintar mediante el comando \type{draw}.

      Esta función se usa añadiendo entre paréntesis un texto:
      \type{textext("...texto...")}. Cuando se procesa, genera un objeto gráfico de
      tipo \type{picture}, que es una imagen en \MetaPost (no un camino o trazo). El
      objeto \type{picture} es una estructura compleja que contiene información
      tipográfica, fuentes incrustadas, colores propios y contornos cerrados
      devueltos por el motor tipográfico de \ConTeXt\.

      \startsubsection[title={Comandos de \ConTeXt\ en el texto}]
    
        En el texto se pueden utilizar comandos estándar de \ConTeXt, como por
        ejemplo \type{framed}, que lo recuadra
    
        \starttyping[option=METAFUN]
          \startMPcode
            fill fullcircle xyscaled (8cm,1cm) withcolor "darkred" ;
            draw textex{"\framed{\bf Texto A}"} withcolor "white" ;
          \stopMPcode
        \stoptyping

        \startMPcode
          fill fullcircle xyscaled (8cm, 1cm) withcolor "green" ;
          draw textext("\bf \framed{This is text A}") withcolor "black" ;
        \stopMPcode

      \stopsubsection
   
    \stopsection
    


%  draw fullcircle xyscaled (4cm, 2cm) ;
%draw fullcircle xscaled 4cm yscaled 2cm ;
%draw fullcircle transformed identity scaled 1 xscaled 4cm yscaled 2cm ; % Interno
%

\stopchapter

\stopcomponent




%%% Local Variables:
%%% coding: utf-8
%%% mode: context
%%% ConTeXt-Mark-version: "IV"
%%% TeX-engine: default
%%% TeX-master: "../ctx-tutorialLuaMetaFun.tex"
%%% End:
